% Canonical integration of agent_logs/native_direct_integral_complete_canonical_20260903.tex.
% The fixed and varying fields remain separately typed; T is a multiplication
% parameter and never a hidden radial derivative.

\subsection{The fixed and varying radial direct integrals}
\label{cusp-native-full-flux-radial-section}

Choose \(T_1>T_0\) so that the complete deep-tail inequalities in
\cref{cusp-native-full-flux-deep-tail-window} hold on
\[
 I:=(T_1,\infty).
\]
Let \(Y=\mathbb R^4/\mathbb Z^4\).  Denote by
\(\mathscr L_\eta\to Y\) the fixed marked Hermitian line bundle defined
by the factor of automorphy in
\cref{cusp-native-full-flux-cocycle}.  The marking diffeomorphism and its
transported bundle isomorphism are
\[
 \Phi_T:Y\longrightarrow F_T,\qquad [y]\longmapsto[P_Ty],
 \qquad
 \widehat\Phi_T:\mathscr L_\eta\longrightarrow\mathscr L_{\eta,T}.
\]
They are fixed by requiring
\[
 \widehat\Phi_T^*\nabla_{\eta,T}
 =\mathrm d-2\pi i(6y^1\,\mathrm dy^4+y^2\,\mathrm dy^3),
\]
so no gauge, metric, lattice, or orientation normalization is hidden in the
identification.  Put
\[
 \mathscr K_T
 :=L^2(F_T,\mathscr L_{\eta,T};
       \mathrm d\operatorname{vol}_{g_T}).
\]

The pullback of the physical volume is
\begin{equation}
 \Phi_T^*(\mathrm d\operatorname{vol}_{g_T})
 =\sqrt{\det(P_T^{\mathsf T}P_T)}\,\mathrm d^4y
 =D_T\,\mathrm d^4y.
\label{native-di-volume-pullback}
\end{equation}
Put
\[
 \mathscr H_*:=L^2(Y,\mathscr L_\eta;\mathrm d^4y),
 \qquad
 \mathscr D_*:=H^2(Y;\mathscr L_\eta),
\]
where the Sobolev space is defined using any one fixed smooth reference
metric and connection on the fixed compact bundle.  Define
\begin{equation}
 U_T:\mathscr K_T\longrightarrow\mathscr H_*,
 \qquad
 U_Ts:=D_T^{1/2}\widehat\Phi_T^*s.
\label{native-di-marking-unitary}
\end{equation}
Then
\[
 \|U_Ts\|_{\mathscr H_*}^2
 =D_T\int_Y|\widehat\Phi_T^*s|^2\,\mathrm d^4y
 =\|s\|_{\mathscr K_T}^2,
\]
and
\[
 U_T^{-1}v
 =D_T^{-1/2}(\widehat\Phi_T^{-1})^*v.
\]
Thus \(U_T\) is an onto unitary, not merely an isometric injection.

Transport the fibre operator to the fixed space:
\begin{equation}
 A_T:=U_T\Delta_{\eta,T}U_T^{-1},
 \qquad
 \operatorname{Dom}(A_T)=\mathscr D_*.
\label{native-di-fixed-fibre}
\end{equation}
Its closed quadratic form is
\begin{equation}
 \widetilde{\mathfrak a}_T[v,w]
 =\int_Y(G_T^{-1})^{jk}
  \langle\nabla_{\eta,j}v,\nabla_{\eta,k}w\rangle\,\mathrm d^4y,
 \qquad G_T=P_T^{\mathsf T}P_T.
\label{native-di-fixed-form}
\end{equation}
Here \(\nabla_{\eta,j}\) is the \(\mathrm dy^j\)-component of the fixed
marked connection displayed above.  The factor \(D_T^{1/2}\) in \(U_T\) is
constant in \(y\), and it cancels
the factor \(D_T\) in \eqref{native-di-volume-pullback}; hence no unrecorded
factor appears in \eqref{native-di-fixed-form}.

\begin{lemma}[Common domain and measurable resolvents]
\label{native-di-common-domain-lemma}
For every compact interval \(K\Subset I\), the family
\[
 T\longmapsto A_T\in\mathcal B(\mathscr D_*,\mathscr H_*)
\]
is norm-continuous, in fact real analytic.  For every
\(z\in\mathbb C\setminus\mathbb R\), the resolvents
\((A_T-z)^{-1}\) are locally norm-continuous as operators on
\(\mathscr H_*\), and hence form a measurable operator field.
\end{lemma}

\begin{proof}
On \(K\), the matrices \(G_T^{-1}\) are uniformly positive definite, and
their entries are real analytic in \(T\).  In fixed local bundle
trivializations, \eqref{native-di-fixed-form} therefore gives a uniformly
elliptic second-order expression whose coefficients depend real
analytically on \(T\).  Elliptic regularity identifies every operator
domain with the same set \(\mathscr D_*\), and the standard local
coordinate estimate for a differential operator of order two gives
norm-continuity in
\(\mathcal B(\mathscr D_*,\mathscr H_*)\).

Uniform ellipticity on \(K\) gives
\[
 \|v\|_{\mathscr D_*}
 \le C_K\bigl(\|A_Tv\|_{\mathscr H_*}
                  +\|v\|_{\mathscr H_*}\bigr),
 \qquad T\in K.
\]
For \(v=(A_T-z)^{-1}f\), self-adjointness gives
\(\|v\|\le|\operatorname{Im}z|^{-1}\|f\|\), while
\(A_Tv=f+zv\).  Consequently
\begin{equation}
 \sup_{T\in K}
 \|(A_T-z)^{-1}\|_{\mathscr H_*\to\mathscr D_*}<\infty.
\label{native-di-elliptic-resolvent-bound}
\end{equation}
The resolvent identity, with its exact sign,
\[
 (A_T-z)^{-1}-(A_R-z)^{-1}
 =-(A_T-z)^{-1}(A_T-A_R)(A_R-z)^{-1},
\]
together with \eqref{native-di-elliptic-resolvent-bound}, proves local
norm-resolvent continuity.  A countable compact exhaustion of \(I\) gives
measurability on all of \(I\).
\end{proof}

We now define the varying measurable Hilbert field rather than presuppose
one.  A section \(T\mapsto s_T\in\mathscr K_T\) is declared measurable
precisely when
\[
 T\longmapsto U_Ts_T\in\mathscr H_*
\]
is strongly measurable.  This is the measurable field pulled back from the
constant separable field \(I\times\mathscr H_*\).  Its direct integral is
\[
 \mathscr H_{\mathrm{var}}
 :=\int_I^\oplus\mathscr K_T\,\mathrm d T.
\]
The fixed radial Hilbert space is
\[
 \mathscr H_{\mathrm{rad}}
 :=L^2(I,\mathrm d T;\mathscr H_*).
\]
By construction,
\begin{equation}
 \mathscr U:=\int_I^\oplus U_T\,\mathrm d T:
 \mathscr H_{\mathrm{var}}\longrightarrow\mathscr H_{\mathrm{rad}},
 \qquad
 (\mathscr Us)(T)=U_Ts_T,
\label{native-di-direct-integral-unitary}
\end{equation}
is an isometry.  It is onto: if
\(\Psi\in\mathscr H_{\mathrm{rad}}\), then
\(s_T=U_T^{-1}\Psi(T)\) is measurable by the very definition of the
varying field and has the same squared norm integral.

On the fixed direct integral, define the maximal multiplication operator
\begin{equation}
 (\widetilde{\mathscr M}_\eta\Psi)(T):=A_T\Psi(T)
\label{native-di-fixed-multiplication}
\end{equation}
with the exact domain
\begin{equation}
\begin{split}
 \operatorname{Dom}(\widetilde{\mathscr M}_\eta)
 :=\{\Psi\in\mathscr H_{\mathrm{rad}}:\;&
 \Psi(T)\in\mathscr D_*\text{ for almost every }T,\\
 &T\mapsto A_T\Psi(T)\text{ is strongly measurable},\\
 &\int_I\|A_T\Psi(T)\|_{\mathscr H_*}^2\,\mathrm d T<\infty\}.
\end{split}
\label{native-di-fixed-maximal-domain}
\end{equation}
On the varying direct integral, define
\begin{equation}
\begin{split}
 \operatorname{Dom}(\mathscr M_\eta)
 :=\{s\in\mathscr H_{\mathrm{var}}:\;&
 s_T\in\operatorname{Dom}(\Delta_{\eta,T})\text{ for almost every }T,\\
 &T\mapsto U_T\Delta_{\eta,T}s_T
       \text{ is strongly measurable},\\
 &\int_I\|\Delta_{\eta,T}s_T\|_{\mathscr K_T}^2\,\mathrm d T<\infty\},
\end{split}
\label{native-di-varying-maximal-domain}
\end{equation}
and
\begin{equation}
 (\mathscr M_\eta s)_T:=\Delta_{\eta,T}s_T
 \quad\text{for almost every }T.
\label{native-di-varying-multiplication}
\end{equation}
Every expression in \eqref{native-di-fixed-multiplication} and
\eqref{native-di-varying-multiplication} is now typed on its actual fibre.
Equations \eqref{native-di-fixed-fibre},
\eqref{native-di-direct-integral-unitary},
\eqref{native-di-fixed-maximal-domain}, and
\eqref{native-di-varying-maximal-domain} give both the domain equality and
the exact intertwining identity
\begin{equation}
 \mathscr U\operatorname{Dom}(\mathscr M_\eta)
 =\operatorname{Dom}(\widetilde{\mathscr M}_\eta),
 \qquad
 \mathscr U\mathscr M_\eta\mathscr U^{-1}
 =\widetilde{\mathscr M}_\eta.
\label{native-di-intertwining}
\end{equation}

\subsubsection{Density and self-adjointness on the maximal domain}

\begin{proposition}[Maximal-domain realization]
\label{native-di-self-adjoint-proposition}
The operators \(\widetilde{\mathscr M}_\eta\) and
\(\mathscr M_\eta\) are densely defined, nonnegative, and self-adjoint on
the domains \eqref{native-di-fixed-maximal-domain} and
\eqref{native-di-varying-maximal-domain}, respectively.
\end{proposition}

\begin{proof}
For \(\chi\in C_c(I)\) and \(v\in\mathscr D_*\), local continuity of
\(T\mapsto A_Tv\) makes it bounded on \(\operatorname{supp}\chi\).
Thus \(\chi v\in\operatorname{Dom}(\widetilde{\mathscr M}_\eta)\).  Since
\(\mathscr D_*\) is dense in \(\mathscr H_*\), the algebraic tensor product
\[
 C_c(I)\mathbin{\odot}\mathscr D_*
\]
is dense in \(L^2(I;\mathscr H_*)\).  Hence the fixed operator is densely
defined.  Fibrewise self-adjointness and Cauchy--Schwarz give, for
\(\Psi,\Theta\) in its domain,
\[
 \langle\widetilde{\mathscr M}_\eta\Psi,\Theta\rangle
 =\int_I\langle A_T\Psi(T),\Theta(T)\rangle\,\mathrm d T
 =\langle\Psi,\widetilde{\mathscr M}_\eta\Theta\rangle,
\]
so it is symmetric.

Fix \(z\in\mathbb C\setminus\mathbb R\).  For
\(F\in\mathscr H_{\mathrm{rad}}\), define pointwise
\begin{equation}
 (\mathcal R_zF)(T):=(A_T-z)^{-1}F(T).
\label{native-di-radial-resolvent}
\end{equation}
The locally norm-continuous resolvent field, approximation of \(F\) by
simple functions, and
\[
 \|(A_T-z)^{-1}\|_{\mathscr H_*\to\mathscr H_*}
 \le|\operatorname{Im}z|^{-1}
\]
show that \(\mathcal R_zF\) is strongly measurable and
\begin{equation}
 \|\mathcal R_zF\|_{\mathscr H_{\mathrm{rad}}}
 \le|\operatorname{Im}z|^{-1}
    \|F\|_{\mathscr H_{\mathrm{rad}}}.
\label{native-di-radial-resolvent-L2}
\end{equation}
Pointwise, \((A_T-z)^{-1}F(T)\in\mathscr D_*\).  The exact identity
\begin{equation}
 A_T(A_T-z)^{-1}=I+z(A_T-z)^{-1}
\label{native-di-maximal-domain-identity}
\end{equation}
gives the global uniform bound
\[
 \|A_T(A_T-z)^{-1}\|
 \le1+\frac{|z|}{|\operatorname{Im}z|}.
\]
The right-hand side of \eqref{native-di-maximal-domain-identity}, applied
to \(F(T)\), is strongly measurable and square-integrable.  Therefore
\(\mathcal R_zF\) belongs to the full maximal domain
\eqref{native-di-fixed-maximal-domain}; this is the domain step that
pointwise inversion alone does not prove.

Fibrewise identities now yield
\[
 (\widetilde{\mathscr M}_\eta-z)\mathcal R_z=I
 \quad\text{on }\mathscr H_{\mathrm{rad}},
 \qquad
 \mathcal R_z(\widetilde{\mathscr M}_\eta-z)=I
 \quad\text{on }\operatorname{Dom}(\widetilde{\mathscr M}_\eta).
\]
In particular, the ranges of
\(\widetilde{\mathscr M}_\eta-i\) and
\(\widetilde{\mathscr M}_\eta+i\) are both all of
\(\mathscr H_{\mathrm{rad}}\).  The range criterion for a densely defined
symmetric operator proves self-adjointness.  Moreover,
\[
 \langle\widetilde{\mathscr M}_\eta\Psi,\Psi\rangle
 =\int_I\langle A_T\Psi(T),\Psi(T)\rangle\,\mathrm d T\ge0,
\]
so it is nonnegative.  The unitary equivalence
\eqref{native-di-intertwining} transfers density, self-adjointness, and
nonnegativity to the varying-fibre operator.
\end{proof}

\subsubsection{Measurable eigenvectors, including at band crossings}

\begin{lemma}[Explicit measurable eigenvector selection]
\label{native-di-measurable-eigenvector-lemma}
For every fixed
\(\mathbf n=(n_+,n_-)\in\mathbb Z_{\ge0}^2\), there is a strongly
measurable map
\[
 I\longrightarrow\mathscr D_*,
 \qquad T\longmapsto\phi_{\mathbf n}(T),
\]
such that
\begin{equation}
 \|\phi_{\mathbf n}(T)\|_{\mathscr H_*}=1,
 \qquad
 A_T\phi_{\mathbf n}(T)
 =E_{\mathbf n}(T)\phi_{\mathbf n}(T)
\label{native-di-measurable-eigenvector}
\end{equation}
for every \(T\in I\).  The construction remains valid when
\(E_{\mathbf n}(T)\) crosses one or more other bands.
\end{lemma}

\begin{proof}
Set
\[
 C_{\mathbf n,T}:=A_T-E_{\mathbf n}(T),
 \qquad \operatorname{Dom}(C_{\mathbf n,T})=\mathscr D_*.
\]
The band \(E_{\mathbf n}\) is continuous.  The proof of
Lemma~\ref{native-di-common-domain-lemma}, applied also to the scalar
shift, shows that
\((C_{\mathbf n,T}-z)^{-1}\) is a measurable resolvent field for every
nonreal \(z\).

Let
\[
 \mathsf P_{\mathbf n}(T)
 :=\mathbf 1_{\{0\}}(C_{\mathbf n,T})
 =\mathbf 1_{\{E_{\mathbf n}(T)\}}(A_T).
\]
This singleton projection is not assumed continuous at a crossing.
Instead, the spectral theorem gives the pointwise strong limit
\begin{equation}
 \mathsf P_{\mathbf n}(T)
 =\operatorname*{s-lim}_{k\to\infty}
  (I+kC_{\mathbf n,T}^2)^{-1},
\label{native-di-kernel-projection}
\end{equation}
and the exact resolvent formula
\begin{equation}
 (I+kC_{\mathbf n,T}^2)^{-1}
 =\frac1{2i\sqrt{k}}
 \left[
 (C_{\mathbf n,T}-i/\sqrt{k})^{-1}
 -(C_{\mathbf n,T}+i/\sqrt{k})^{-1}
 \right]
\label{native-di-kernel-projection-resolvent}
\end{equation}
shows that every approximating field is strongly measurable.  Therefore
\(T\mapsto\mathsf P_{\mathbf n}(T)v\) is strongly measurable for every
\(v\in\mathscr H_*\).  Proposition
\ref{cusp-native-full-flux-landau-theorem} says that
\(\mathsf P_{\mathbf n}(T)\ne0\) for every \(T\), including crossings.

Choose a countable dense sequence \((v_k)_{k\ge1}\) in the separable
space \(\mathscr H_*\), and define the measurable partition
\[
 X_k=
 \{T:\|\mathsf P_{\mathbf n}(T)v_k\|>0\}
 \cap
 \bigcap_{\ell<k}
 \{T:\|\mathsf P_{\mathbf n}(T)v_\ell\|=0\}.
\]
The sets \(X_k\) are pairwise disjoint and cover \(I\): a nonzero bounded
projection cannot annihilate a dense set.  On \(X_k\), set
\[
 \phi_{\mathbf n}(T)
 =\frac{\mathsf P_{\mathbf n}(T)v_k}
 {\|\mathsf P_{\mathbf n}(T)v_k\|}.
\]
Countable measurable patching gives a strongly measurable unit vector.
The range of the singleton spectral projection lies in
\(\operatorname{Dom}(A_T)\) and satisfies
\((A_T-E_{\mathbf n}(T))\phi_{\mathbf n}(T)=0\), which proves
\eqref{native-di-measurable-eigenvector}.
\end{proof}

\subsubsection{Full spectrum and spectral type}

\begin{theorem}[Purely continuous full half-line spectrum]
\label{native-di-main-theorem}
For both the fixed-fibre and varying-fibre maximal radial multiplication
operators,
\begin{equation}
\begin{aligned}
 \sigma(\widetilde{\mathscr M}_\eta)
 &=\sigma_{\mathrm{ess}}(\widetilde{\mathscr M}_\eta)
 =\sigma_{\mathrm c}(\widetilde{\mathscr M}_\eta)
 =[0,\infty),\\
 \sigma_{\mathrm p}(\widetilde{\mathscr M}_\eta)
 &=\varnothing,
 \qquad
 \ker\widetilde{\mathscr M}_\eta=\{0\},
\end{aligned}
\label{native-di-fixed-spectrum-result}
\end{equation}
and
\begin{equation}
\begin{aligned}
 \sigma(\mathscr M_\eta)
 &=\sigma_{\mathrm{ess}}(\mathscr M_\eta)
 =\sigma_{\mathrm c}(\mathscr M_\eta)
 =[0,\infty),\\
 \sigma_{\mathrm p}(\mathscr M_\eta)
 &=\varnothing,
 \qquad
 \ker\mathscr M_\eta=\{0\}.
\end{aligned}
\label{native-di-varying-spectrum-result}
\end{equation}
Here \(\sigma_{\mathrm{ess}}\) is the Weyl essential spectrum.
\end{theorem}

\begin{proof}
It is enough to prove \eqref{native-di-fixed-spectrum-result}, because
\eqref{native-di-intertwining} is a unitary equivalence.
Proposition~\ref{native-di-self-adjoint-proposition} and nonnegativity give
\[
 \sigma(\widetilde{\mathscr M}_\eta)\subset[0,\infty).
\]

\paragraph{The zero-energy singular sequence.}
Let \(a_j=T_1+2j+1\).  Using the measurable normalized ground eigenvector
from Lemma~\ref{native-di-measurable-eigenvector-lemma}, define
\[
 \Psi_j(T)
 :=\mathbf 1_{[a_j,a_j+1]}(T)\phi_{(0,0)}(T).
\]
These vectors have norm one and have pairwise disjoint supports, hence are
orthonormal and converge weakly to zero.  They belong to the maximal
domain because
\[
 A_T\Psi_j(T)=
 \mathbf 1_{[a_j,a_j+1]}(T)E_0(T)\phi_{(0,0)}(T)
\]
is strongly measurable and square-integrable.  From
\eqref{cusp-native-full-flux-ground-energy},
\[
 \|\widetilde{\mathscr M}_\eta\Psi_j\|^2
 =\int_{a_j}^{a_j+1}E_0(T)^2\,\mathrm d T
 =\frac{196\pi^2}{a_j^2}+O(a_j^{-3})
 \longrightarrow0.
\]
Thus \((\Psi_j)\) is a singular Weyl sequence at zero, and
\[
 0\in\sigma_{\mathrm{ess}}(\widetilde{\mathscr M}_\eta).
\]

\paragraph{Positive energy: escaping crossings before disjoint supports.}
Fix \(\lambda>0\), and fix once and for all
\(T_\bullet=T_1+1\in I\).  For \(n\in\mathbb Z_{\ge0}\), retain the
specific canonical band
\begin{equation}
 e_n(T):=E_{(0,n)}(T)
 =\mathcal B_{T,+}+(2n+1)\mathcal B_{T,-}.
\label{native-di-crossing-band}
\end{equation}
Since \(\mathcal B_{T_\bullet,-}>0\),
\[
 e_n(T_\bullet)\longrightarrow\infty
 \quad(n\to\infty).
\]
For each fixed \(n\), equations
\eqref{cusp-native-full-flux-Bminus-expansion} and
\eqref{cusp-native-full-flux-Bplus-expansion} give
\[
 e_n(T)\longrightarrow0
 \quad(T\to\infty).
\]
Hence, for every sufficiently large \(n\), the intermediate value theorem
gives at least one crossing parameter \(r_n>T_\bullet\) satisfying
\begin{equation}
 e_n(r_n)=\lambda.
\label{native-di-crossing}
\end{equation}

We prove explicitly that every such choice of crossing parameters escapes
to infinity.  Fix \(M>T_\bullet\).  Positivity and continuity give
\[
 c_M:=\min_{T\in[T_\bullet,M]}\mathcal B_{T,-}>0.
\]
For every \(T\in[T_\bullet,M]\),
\[
 e_n(T)\ge(2n+1)c_M.
\]
Thus, once \((2n+1)c_M>\lambda\), equation
\eqref{native-di-crossing} has no solution in
\([T_\bullet,M]\), and the chosen \(r_n\) must satisfy \(r_n>M\).
Since \(M\) was arbitrary,
\begin{equation}
 r_n\longrightarrow\infty.
\label{native-di-crossings-escape}
\end{equation}

Only now, before choosing any support, pass to a subsequence
\(n_j\to\infty\) such that, with \(r_j:=r_{n_j}\),
\begin{equation}
 r_{j+1}>r_j+2.
\label{native-di-separated-crossings}
\end{equation}
Continuity of \(e_{n_j}\) at \(r_j\) permits the subsequent choice of
\(\varepsilon_j\in(0,\tfrac14)\) such that
\[
 J_j:=(r_j-\varepsilon_j,r_j+\varepsilon_j)\Subset I,
 \qquad
 \sup_{T\in J_j}|e_{n_j}(T)-\lambda|<\frac1j.
\]
The order of construction matters: the separation
\eqref{native-di-separated-crossings} makes the intervals \(J_j\)
pairwise disjoint.

Take the measurable normalized eigenvector
\(\phi_{(0,n_j)}(T)\) from
Lemma~\ref{native-di-measurable-eigenvector-lemma}; the singleton
projection used there remains valid if the selected band crosses another
one inside \(J_j\).  Define
\begin{equation}
 \Theta_j(T)
 :=|J_j|^{-1/2}\mathbf 1_{J_j}(T)\phi_{(0,n_j)}(T).
\label{native-di-positive-Weyl-sequence}
\end{equation}
The sections \(\Theta_j\) have norm one and are orthogonal by disjointness
of their supports.  They lie in the maximal domain because
\[
 A_T\Theta_j(T)
 =|J_j|^{-1/2}\mathbf 1_{J_j}(T)
   e_{n_j}(T)\phi_{(0,n_j)}(T)
\]
is strongly measurable and square-integrable.  Moreover,
\[
 \|(\widetilde{\mathscr M}_\eta-\lambda)\Theta_j\|^2
 =\frac1{|J_j|}\int_{J_j}
   |e_{n_j}(T)-\lambda|^2\,\mathrm d T
 <\frac1{j^2}.
\]
Every orthonormal sequence is weakly null, so
\((\Theta_j)\) is a singular Weyl sequence at \(\lambda\).
Thus every \(\lambda>0\) belongs to the Weyl essential spectrum.  Together
with the zero sequence and nonnegativity, this proves
\begin{equation}
 \sigma(\widetilde{\mathscr M}_\eta)
 =\sigma_{\mathrm{ess}}(\widetilde{\mathscr M}_\eta)
 =[0,\infty).
\label{native-di-spectrum-essential}
\end{equation}

\paragraph{Real analyticity, absence of flat bands, and no eigenvalues.}
Equation \eqref{cusp-native-full-flux-Xi-identity} writes
\[
 \Xi_T=\sqrt{S_T^2+24D_T}
\]
with a strictly positive real-analytic radicand.  The positive square-root
function is real analytic on \((0,\infty)\).  Hence
\(\Xi_T\), \(\mathcal B_{T,+}\), \(\mathcal B_{T,-}\), and every
\(E_{\mathbf n}(T)\) are real analytic on the connected interval \(I\).
For each fixed \(\mathbf n\),
\[
 E_{\mathbf n}(T)>0\quad(T<\infty),
 \qquad
 \lim_{T\to\infty}E_{\mathbf n}(T)=0.
\]
It follows that no band is constant: a constant band would have constant
value equal to its limit zero, contradicting strict positivity at every
finite \(T\).  This proves the absence of flat bands, rather than assuming
it.

Fix \(\lambda\in\mathbb R\).  For every \(\mathbf n\), the real-analytic
function \(E_{\mathbf n}-\lambda\) is not identically zero, so its zero set
in \(I\) is discrete and has Lebesgue measure zero.  Since
\(\mathbb Z_{\ge0}^2\) is countable,
\begin{equation}
 Z_\lambda
 :=\{T\in I:\lambda\in\sigma_{\mathrm p}(A_T)\}
 =\bigcup_{\mathbf n\in\mathbb Z_{\ge0}^2}
   \{T:E_{\mathbf n}(T)=\lambda\}
\label{native-di-null-fibre-level-set}
\end{equation}
is a null set.

If
\(\widetilde{\mathscr M}_\eta\Psi=\lambda\Psi\), then
\((A_T-\lambda)\Psi(T)=0\) for almost every \(T\).  At every point where
\(\Psi(T)\ne0\), Proposition
\ref{cusp-native-full-flux-landau-theorem} forces \(T\in Z_\lambda\).
Thus the essential support of \(\Psi\) is contained in the null set
\eqref{native-di-null-fibre-level-set}, and \(\Psi=0\).  A self-adjoint
operator has only real eigenvalues, so
\[
 \sigma_{\mathrm p}(\widetilde{\mathscr M}_\eta)=\varnothing,
 \qquad
 \ker\widetilde{\mathscr M}_\eta=\{0\}.
\]
Self-adjoint operators have empty residual spectrum.  Combining this with
\eqref{native-di-spectrum-essential} proves that the entire spectrum is
continuous and gives \eqref{native-di-fixed-spectrum-result}.  Unitary
equivalence gives \eqref{native-di-varying-spectrum-result}.
\end{proof}

\subsubsection{Exact interpretive boundary}

\begin{remark}[No radial derivative and no interacting Yang--Mills claim]
The variable \(T\) in
\eqref{native-di-fixed-multiplication} and
\eqref{native-di-varying-multiplication} is a direct-integral parameter.
Neither operator contains a \(T\)-derivative.  In particular, there is no
term such as \(-\partial_T^2\), no radial Sobolev condition, no cusp
boundary condition, and no interaction between distinct \(T\)-fibres.
This is why characteristic functions of radial intervals legitimately
occur in the maximal-domain Weyl sequences above.

The theorem is therefore a complete spectral statement for the specified
auxiliary decomposable family of native full-flux magnetic Laplacians.  It
does not construct, identify, or determine the spectrum or mass gap of an
interacting quantum Yang--Mills Hamiltonian, and it makes no
Osterwalder--Schrader reconstruction claim.
\end{remark}
